\chapter{Geometría analítica}\label{ch:g10:coordgeom}

La gran idea de Descartes fue describir los puntos mediante parejas de
números, convirtiendo los problemas de geometría en cálculos. Con solo dos
fórmulas, la del \omterm{prop:g10:coordgeom:midpoint}{punto medio} y la de la distancia, se puede demostrar que
un triángulo es isósceles, que un cuadrilátero es un paralelogramo o que
tres puntos están alineados, sin trazar ni una sola línea auxiliar.

\section{Coordenadas en el plano}

\begin{definition}[Sistema de coordenadas]\label{def:g10:coordgeom:system}
Un \emph{sistema de coordenadas}\index{sistema de coordenadas} del plano
consta de un origen $O$ y dos ejes graduados que pasan por $O$: el eje
horizontal, o eje $x$, y el eje vertical, o eje $y$. Cada punto $M$ tiene
entonces una única pareja de \emph{coordenadas}\index{coordenadas}
$(x, y)$: su \emph{abscisa}\index{abscisa} $x$ y su
\emph{ordenada}\index{ordenada} $y$. El sistema es
\emph{ortonormal}\index{sistema ortonormal} cuando los ejes son
perpendiculares y llevan la misma unidad de longitud. Todos los sistemas
de este capítulo son ortonormales.
\end{definition}

\begin{omfigure}
\begin{tikzpicture}[scale=0.85]
  \draw[thin, gray!40] (-2.4,-1.4) grid (4.4,3.4);
  \draw[-Stealth] (-2.6,0) -- (4.6,0) node[below, font=\small] {$x$};
  \draw[-Stealth] (0,-1.6) -- (0,3.6) node[left, font=\small] {$y$};
  \node[below left, font=\small] at (0,0) {$O$};
  \fill[omDef] (3,2) circle (2pt);
  \node[above right, font=\small, omDef] at (3,2) {$M(3, 2)$};
  \draw[dashed, omDef] (3,0) -- (3,2) -- (0,2);
  \fill[omThm] (-2,-1) circle (2pt);
  \node[below left, font=\small, omThm] at (-2,-1) {$N(-2, -1)$};
  \foreach \x in {1,2,3,4} \node[below, font=\footnotesize] at (\x,0) {\x};
  \foreach \y in {1,2,3} \node[left, font=\footnotesize] at (0,\y) {\y};
\end{tikzpicture}

{\small Cada punto del plano queda localizado por dos números: primero la
\omterm{def:g10:coordgeom:system}{abscisa} (horizontal) y después la \omterm{def:g10:coordgeom:system}{ordenada} (vertical).}
\end{omfigure}

\section{Punto medio de un segmento}

\begin{proposition}[Fórmula del punto medio]\label{prop:g10:coordgeom:midpoint}
Sean $A(x_A, y_A)$ y $B(x_B, y_B)$. El \emph{punto medio}\index{punto
medio} $I$ del segmento $[AB]$ tiene \omterm{def:g10:coordgeom:system}{coordenadas}
\[
I\left(\frac{x_A + x_B}{2},\ \frac{y_A + y_B}{2}\right).
\]
\end{proposition}

\begin{proof}
Fijémonos en las \omterm{def:g10:coordgeom:system}{coordenadas} horizontales. Al ir de $A$ a $B$, la \omterm{def:g10:coordgeom:system}{abscisa}
varía en $x_B - x_A$; el punto situado a mitad de camino tiene \omterm{def:g10:coordgeom:system}{abscisa}
\[
x_A + \frac{x_B - x_A}{2} = \frac{2x_A + x_B - x_A}{2}
= \frac{x_A + x_B}{2},
\]
la media de las dos \omterm{def:g10:coordgeom:system}{abscisas}. El mismo cálculo vale para las \omterm{def:g10:coordgeom:system}{ordenadas}.
\end{proof}

\begin{example}\label{ex:g10:coordgeom:midpoint}
Con $A(-1, 4)$ y $B(5, -2)$, el \omterm{prop:g10:coordgeom:midpoint}{punto medio} de $[AB]$ es
\[
I\left(\frac{-1 + 5}{2},\ \frac{4 + (-2)}{2}\right) = I(2, 1).
\]
La fórmula también funciona al revés: si $A(-1, 4)$ y el \omterm{prop:g10:coordgeom:midpoint}{punto medio} es
$I(2, 1)$, entonces $B$ cumple $\frac{-1 + x_B}{2} = 2$ y
$\frac{4 + y_B}{2} = 1$, luego $x_B = 5$ e $y_B = -2$: el punto $B$ es el
\emph{simétrico} de $A$ respecto de $I$.
\end{example}

\section{Distancia entre dos puntos}

\begin{theorem}[Fórmula de la distancia]\label{thm:g10:coordgeom:distance}
En un \omterm{def:g10:coordgeom:system}{sistema ortonormal}, la distancia entre $A(x_A, y_A)$ y
$B(x_B, y_B)$ es
\[
AB = \sqrt{(x_B - x_A)^2 + (y_B - y_A)^2}.
\]
\end{theorem}

\begin{proof}
Sea $C$ el punto $(x_B, y_A)$: está a la misma altura que $A$ y tiene la
misma \omterm{def:g10:coordgeom:system}{abscisa} que $B$, así que el triángulo $ACB$ tiene un ángulo recto en
$C$. Sus catetos son segmentos horizontales y verticales, de longitudes
$AC = \abs{x_B - x_A}$ y $CB = \abs{y_B - y_A}$. Por el teorema de
Pitágoras,
\[
AB^2 = AC^2 + CB^2 = (x_B - x_A)^2 + (y_B - y_A)^2,
\]
y al tomar la raíz cuadrada (los dos miembros son no negativos) se obtiene
la fórmula.
\end{proof}

\begin{omfigure}
\begin{tikzpicture}[scale=0.85]
  \draw[thin, gray!40] (-0.4,-0.4) grid (5.4,3.4);
  \draw[-Stealth] (-0.6,0) -- (5.6,0) node[below, font=\small] {$x$};
  \draw[-Stealth] (0,-0.6) -- (0,3.6) node[left, font=\small] {$y$};
  \coordinate (A) at (1,1);
  \coordinate (B) at (5,3);
  \coordinate (C) at (5,1);
  \draw[omDef, thick] (A) -- (B);
  \draw[omThm, thick] (A) -- (C) node[midway, below, font=\small]
    {$\abs{x_B - x_A}$};
  \draw[omThm, thick] (C) -- (B) node[midway, right, font=\small]
    {$\abs{y_B - y_A}$};
  \draw[thin] (C) ++(-0.25,0) -- ++(0,0.25) -- ++(0.25,0);
  \fill (A) circle (2pt) node[above left, font=\small] {$A$};
  \fill (B) circle (2pt) node[above, font=\small] {$B$};
  \fill (C) circle (2pt) node[below right, font=\small] {$C$};
\end{tikzpicture}

{\small La fórmula de la distancia es el teorema de Pitágoras aplicado al
triángulo rectángulo $ACB$, construido sobre un cateto horizontal y otro
vertical.}
\end{omfigure}

\begin{example}\label{ex:g10:coordgeom:distance}
Con $A(1, 1)$ y $B(5, 3)$:
\[
AB = \sqrt{(5-1)^2 + (3-1)^2} = \sqrt{16 + 4} = \sqrt{20} = 2\sqrt 5 .
\]
Las distancias suelen dejarse en forma exacta (con la raíz); redondea solo
al final, si hace falta una respuesta decimal.
\end{example}

\begin{remark}
El orden de los puntos no importa: $(x_A - x_B)^2 = (x_B - x_A)^2$. ¡Pero
no te olvides de los cuadrados! La distancia \emph{no} es
$\abs{x_B - x_A} + \abs{y_B - y_A}$.
\end{remark}

\section{Las fórmulas al servicio de la geometría}

\begin{method}[Naturaleza de un triángulo]\label{met:g10:coordgeom:triangle}
Dados tres puntos $A$, $B$, $C$ por sus \omterm{def:g10:coordgeom:system}{coordenadas}:
\begin{enumerate}
  \item calcula los tres cuadrados de las longitudes $AB^2$, $AC^2$,
        $BC^2$ con la fórmula de la distancia (dejarlo en cuadrados evita
        las raíces);
  \item dos cuadrados iguales $\Rightarrow$ el triángulo es
        \emph{isósceles};
  \item si el mayor de los cuadrados es la suma de los otros dos, el
        triángulo es \emph{rectángulo} en el vértice opuesto al lado más
        largo, por el recíproco del teorema de Pitágoras.
\end{enumerate}
\end{method}

\begin{example}\label{ex:g10:coordgeom:triangle}
Sean $A(0, 1)$, $B(4, 3)$ y $C(2, -3)$. Entonces
\begin{align*}
AB^2 &= (4-0)^2 + (3-1)^2 = 16 + 4 = 20, \\
AC^2 &= (2-0)^2 + (-3-1)^2 = 4 + 16 = 20, \\
BC^2 &= (2-4)^2 + (-3-3)^2 = 4 + 36 = 40 .
\end{align*}
En primer lugar, $AB^2 = AC^2$, luego $AB = AC$: el triángulo es isósceles
en $A$. Además $AB^2 + AC^2 = 40 = BC^2$, así que por el recíproco del
teorema de Pitágoras también es rectángulo en $A$.
\end{example}

\begin{method}[Reconocer un paralelogramo]\label{met:g10:coordgeom:parallelogram}
Un cuadrilátero $ABCD$ es un paralelogramo exactamente cuando sus
diagonales $[AC]$ y $[BD]$ tienen el mismo \omterm{prop:g10:coordgeom:midpoint}{punto medio}. Así pues: calcula
los dos \omterm{prop:g10:coordgeom:midpoint}{puntos medios} con la fórmula correspondiente y compáralos.
\end{method}

\begin{example}\label{ex:g10:coordgeom:parallelogram}
Sean $A(-1, 0)$, $B(2, 2)$, $C(5, 1)$ y $D(2, -1)$. El \omterm{prop:g10:coordgeom:midpoint}{punto medio} de
$[AC]$ es $\left(\frac{-1+5}{2}, \frac{0+1}{2}\right) = (2, \frac12)$; el
de $[BD]$ es $\left(\frac{2+2}{2}, \frac{2-1}{2}\right) =
(2, \frac12)$. Coinciden, así que $ABCD$ es un paralelogramo.
\end{example}

\begin{omfigure}
\begin{tikzpicture}[scale=0.75]
  \draw[thin, gray!40] (-1.4,-1.4) grid (5.4,2.4);
  \draw[-Stealth] (-1.6,0) -- (5.6,0);
  \draw[-Stealth] (0,-1.6) -- (0,2.6);
  \coordinate (A) at (-1,0);
  \coordinate (B) at (2,2);
  \coordinate (C) at (5,1);
  \coordinate (D) at (2,-1);
  \draw[omDef, thick] (A) -- (B) -- (C) -- (D) -- cycle;
  \draw[omThm, dashed] (A) -- (C);
  \draw[omThm, dashed] (B) -- (D);
  \fill (A) circle (2pt) node[left, font=\small] {$A$};
  \fill (B) circle (2pt) node[above, font=\small] {$B$};
  \fill (C) circle (2pt) node[right, font=\small] {$C$};
  \fill (D) circle (2pt) node[below, font=\small] {$D$};
  \fill[omThm] (2,0.5) circle (2pt) node[above right, font=\small] {$I$};
\end{tikzpicture}

{\small $ABCD$ es un paralelogramo porque sus dos diagonales comparten el
\omterm{prop:g10:coordgeom:midpoint}{punto medio} $I(2, \frac12)$.}
\end{omfigure}

\section{Ejercicios}

\begin{exercise}[$\star$]\label{exo:g10:coordgeom:1}
Sean $A(2, 5)$ y $B(-4, 1)$. Calcula las \omterm{def:g10:coordgeom:system}{coordenadas} del \omterm{prop:g10:coordgeom:midpoint}{punto medio} de
$[AB]$ y la distancia $AB$.
\end{exercise}

\begin{exercise}[$\star$]\label{exo:g10:coordgeom:2}
Sean $A(3, -2)$ e $I(1, 2)$. Halla las \omterm{def:g10:coordgeom:system}{coordenadas} del punto $B$ tal que
$I$ sea el \omterm{prop:g10:coordgeom:midpoint}{punto medio} de $[AB]$.
\end{exercise}

\begin{exercise}[$\star$]\label{exo:g10:coordgeom:3}
¿Cuál de los puntos $P(4, 1)$, $Q(-3, 2)$ y $R(0, -5)$ está más cerca del
origen $O(0,0)$? Responde con distancias exactas.
\end{exercise}

\begin{exercise}[$\star$]\label{exo:g10:coordgeom:4}
Representa los puntos $A(1, 2)$, $B(5, 4)$, $C(7, 0)$ en papel
cuadriculado y demuestra después con un cálculo que el triángulo $ABC$ es
isósceles. ¿Es rectángulo?
\end{exercise}

\begin{exercise}[$\star\star$]\label{exo:g10:coordgeom:5}
Sean $A(-2, 1)$, $B(1, 3)$, $C(4, 1)$ y $D(1, -1)$.
\begin{enumerate}
  \item Demuestra que $ABCD$ es un paralelogramo.
  \item Calcula $AB$ y $BC$. ¿Es $ABCD$ un rombo (todos los lados
        iguales)?
  \item Calcula $AC$ y $BD$. ¿Es $ABCD$ un rectángulo (diagonales
        iguales)?
\end{enumerate}
\end{exercise}

\begin{exercise}[$\star\star$]\label{exo:g10:coordgeom:6}
La circunferencia de centro $\Omega(2, 1)$ y radio $5$ está formada por
todos los puntos que distan $5$ de $\Omega$. ¿Cuáles de los puntos
$A(5, 5)$, $B(-1, -3)$, $C(6, 2)$ están sobre ella? ¿Cuáles en su
interior? ¿Cuáles en su exterior?
\end{exercise}

\begin{exercise}[$\star\star$]\label{exo:g10:coordgeom:7}
Sean $A(1, 1)$ y $B(7, 5)$. Halla todos los puntos $M(x, 0)$ del eje $x$
que equidistan de $A$ y de $B$. (Escribe $MA^2 = MB^2$ y despeja $x$.)
\end{exercise}

\begin{exercise}[$\star\star$]\label{exo:g10:coordgeom:8}
Sean $A(0, 3)$ y $B(4, 1)$.
\begin{enumerate}
  \item Calcula las \omterm{def:g10:coordgeom:system}{coordenadas} del \omterm{prop:g10:coordgeom:midpoint}{punto medio} $I$ de $[AB]$.
  \item Comprueba que $M(2, 2) = I$ y verifica después, calculando $MA$ y
        $MB$, que $M$ equidista de $A$ y de $B$.
  \item Halla un segundo punto, sobre el eje $y$, que equidiste de $A$ y
        de $B$.
\end{enumerate}
\end{exercise}

\begin{exercise}[$\star\star$]\label{exo:g10:coordgeom:9}
Se dan los puntos $A(-1, -1)$, $B(3, 1)$ y $C(11, 5)$. Calcula $AB$, $BC$
y $AC$ y deduce que $A$, $B$, $C$ están alineados. (Tres puntos están
alineados exactamente cuando la mayor de las tres distancias es la suma de
las otras dos.)
\end{exercise}

\begin{exercise}[$\star\star\star$]\label{exo:g10:coordgeom:10}
Sean $A(a, 0)$ y $B(0, b)$ con $a, b > 0$, y sea $I$ el \omterm{prop:g10:coordgeom:midpoint}{punto medio} de
$[AB]$. Demuestra con un cálculo que $OI = IA = IB$, donde $O$ es el
origen. (Esto demuestra un teorema clásico: en un triángulo rectángulo, el
\omterm{prop:g10:coordgeom:midpoint}{punto medio} de la hipotenusa equidista de los tres vértices.)
\end{exercise}

\section{Problema: El puente de Descartes --- ecuaciones de
circunferencias y cómo saber dónde estás}

\begin{problem}[{Problema de fin de semana --- la circunferencia se
convierte en una ecuación, los viejos teoremas en cálculos y tres balizas
localizan un punto: trilateración por álgebra}]
\label{pb:g10:coordgeom:1}
En 1637, René Descartes tendió un puente entre dos continentes: toda curva
de la geometría pasó a ser una \emph{\omterm{def:g10:algebra:equation}{ecuación}} y toda cuestión geométrica,
un cálculo. Este problema recorre el puente en los dos sentidos: deduce la
\omterm{def:g10:algebra:equation}{ecuación} de la circunferencia a partir de la fórmula de la distancia
(\cref{thm:g10:coordgeom:distance}), vuelve a demostrar teoremas clásicos
en tres líneas de álgebra y termina donde el puente soporta hoy más
tráfico: calcular una posición a partir de señales de distancia, el
corazón, en versión plana, del GPS.

\medskip
\noindent\textbf{Parte I --- La circunferencia recibe una \omterm{def:g10:algebra:equation}{ecuación}.}
\begin{enumerate}
  \item Un punto $M(x, y)$ está en la circunferencia de centro
        $\Omega(2, 1)$ y radio $5$ exactamente cuando $\Omega M = 5$.
        Eleva al cuadrado esa condición para obtener la \omterm{def:g10:algebra:equation}{ecuación} de la
        circunferencia:
        \[
        (x - 2)^2 + (y - 1)^2 = 25 .
        \]
  \item Comprueba la pertenencia: ¿cuáles de los puntos $A(5, 5)$,
        $B(6, 4)$, $D(4, 4)$ están sobre la circunferencia, cuáles dentro
        y cuáles fuera?
  \item Una circunferencia disfrazada: completa los cuadrados
        (\cref{pb:g10:algebra:1}) en
        \[
        x^2 + y^2 - 6x + 4y - 12 = 0
        \]
        y da su centro y su radio.
  \item Haz lo mismo con $x^2 + y^2 - 6x + 4y + 14 = 0$. ¿Qué revela la
        forma completada? Enuncia el criterio: ¿cuándo describe
        $(x - h)^2 + (y - k)^2 = c$ una circunferencia, cuándo un único
        punto y cuándo nada en absoluto?
  \item Corta la circunferencia de la pregunta 3 con la recta horizontal
        $y = 3$: sustituye y resuelve. ¿Cuántos puntos salen y qué nombre
        geométrico recibe una recta así?
\end{enumerate}

\medskip
\noindent\textbf{Parte II --- Viejos teoremas en tres líneas.}
\begin{enumerate}[resume]
  \item El teorema de la circunferencia del volumen anterior (un ángulo
        recto en toda semicircunferencia), demostrado de nuevo por
        álgebra: sean $A(-r, 0)$ y $B(r, 0)$ los extremos de un diámetro y
        $M(x, y)$ un punto cualquiera de la circunferencia
        $x^2 + y^2 = r^2$. Calcula $MA^2 + MB^2$ y concluye con el
        recíproco de Pitágoras que $\widehat{AMB}$ es recto.
  \item El teorema de la mediana: con $A(-a, 0)$ y $B(a, 0)$ (de modo que
        el \omterm{prop:g10:coordgeom:midpoint}{punto medio} de $[AB]$ es el origen $I$), demuestra que para
        todo punto $M(x, y)$:
        \[
        MA^2 + MB^2 = 2\,MI^2 + \frac{AB^2}{2} .
        \]
  \item Deduce el lugar geométrico de los puntos $M$ con
        $MA^2 + MB^2 = 26$ cuando $AB = 6$: ¿qué curva, con qué centro y
        qué radio?
  \item Otro lugar geométrico: $A(0,0)$, $B(3,0)$; halla todos los puntos
        con $MA = 2\,MB$. (Eleva al cuadrado, desarrolla, completa los
        cuadrados: aparece una circunferencia famosa, que Apolonio ya
        conocía sin \omterm{def:g10:coordgeom:system}{coordenadas}.)
  \item En una o dos frases: ¿qué cambia el puente de Descartes en la
        \emph{manera} de demostrar teoremas? Compara la pregunta 6 con la
        demostración mediante un rectángulo del teorema de la
        circunferencia del volumen anterior (un ángulo recto en toda
        semicircunferencia).
\end{enumerate}

\medskip
\noindent\textbf{Parte III --- Tres balizas te encuentran.}
Tu posición $M(x, y)$ es desconocida. La baliza $P(0, 0)$ mide tu
distancia y obtiene $5$; la baliza $Q(6, 0)$ obtiene también $5$.
\begin{enumerate}[resume]
  \item Escribe las dos ecuaciones de circunferencia, réstalas y observa
        cómo se cancelan los cuadrados: ¿qué \omterm{def:g10:algebra:equation}{ecuación} sencilla sobrevive?
        Combínala con una de las circunferencias para hallar las
        \emph{dos} posiciones candidatas.
  \item Una tercera baliza, $R(0, 8)$, mide tu distancia y obtiene $5$.
        Calcula su distancia a cada candidata y decide dónde estás.
  \item Explica por qué restar dos ecuaciones de circunferencia da
        \emph{siempre} la \omterm{def:g10:algebra:equation}{ecuación} de una recta (¿qué términos se
        cancelan?) y qué recta es esa, geométricamente, cuando las dos
        circunferencias se cortan en dos puntos.
  \item El posicionamiento por satélite real funciona en el espacio y con
        relojes imperfectos: cada satélite aporta (a través del tiempo de
        viaje de la señal) una distancia y, por tanto, una esfera.
        ¿Cuántas incógnitas tiene un receptor (la posición \emph{y} el
        error de su propio reloj) y por qué el GPS escucha entonces al
        menos a \emph{cuatro} satélites?
  \item Precisión: supón que la distancia de la baliza $P$ es realmente
        $5.1$ en lugar de $5$ (todo lo demás igual). Rehaz la resta de la
        pregunta 11 para hallar el nuevo $x$ y cuantifica cuánto se ha
        desplazado la estimación. (Compara con las ideas sobre el
        presupuesto de errores del \cref{pb:g10:numbers:1}.)
\end{enumerate}

\medskip
\noindent\textbf{Parte IV --- La caja de herramientas del topógrafo.}
\begin{enumerate}[resume]
  \item El camino más corto para ir a beber: un campamento en $A(1, 5)$,
        un establo en $B(7, 3)$ y un río recto a lo largo del eje $y = 0$.
        Para ir de $A$ al río y después a $B$ andando lo menos posible:
        refleja $B$ respecto del río en $B'$ y halla dónde corta el
        segmento $[AB']$ al río. Da el punto de corte y la longitud total
        mínima. (La idea es la del ejercicio del rebote de billar del
        volumen anterior, ahora totalmente calculable.)
  \item Confirma que es mínima frente a una rival: calcula el camino total
        pasando por $Q(3, 0)$ y comprueba que pierde frente a tu
        respuesta.
  \item Clasifica por completo el cuadrilátero $A(1,1)$, $B(4,2)$,
        $C(5,5)$, $D(2,4)$: ¿paralelogramo?, ¿rombo?, ¿rectángulo?,
        ¿cuadrado? (\cref{met:g10:coordgeom:parallelogram},
        \cref{met:g10:coordgeom:triangle}; compara también las
        diagonales.)
  \item Dados $A(-1, 2)$, $B(3, 4)$, $C(6, 0)$, halla $D$ tal que $ABCD$
        sea un paralelogramo (el truco del \omterm{prop:g10:coordgeom:midpoint}{punto medio} de las diagonales
        del problema de fin de semana del volumen anterior sobre los medios
        giros, ahora en fórmulas).
  \item Final: tres herramientas, la fórmula del \omterm{prop:g10:coordgeom:midpoint}{punto medio}, la de la
        distancia y la resta de ecuaciones de circunferencia. Di, para
        cada una, qué tipo de pregunta ha resuelto en este problema y qué
        añaden a la caja los dos capítulos siguientes (direcciones y
        \omterm{def:g10:reffunc:affine}{pendientes}: los vectores y las ecuaciones de rectas).
\end{enumerate}
\end{problem}
